% ===================================================================
%  TABELO N-RO 6 — La interna domo. — La meblaro. — La nutrado. —
%  La lumigado.
%  Traduction 2026 d'apres texto/io, controlee sur texto/fr et
%  texto/es. Consigne : voir texto/eo/10-tabelo-01.tex.
%
%  Deux notes, toutes deux marquees « (*) », et deux appels : celui de
%  « babo » au bloc c1-12-1, celui de « lambrequino » au bloc c2-01-1.
%  L'ordre les departage, comme dans les autres colonnes.
%
%  LA FEMME DE CHAMBRE N'A PAS DE GROS PLAN, et c'est normal. L'ido du
%  bloc c1-06-1 porte « la chambristino (6) », le francais « la femme
%  de chambre (16) » : c'est l'ido qui se trompe, son propre numero 6
%  etant le savon du bloc c1-02-1. L'esperanto suit donc le francais
%  et ecrit (16) ; la planche ne porte pas de 16 releve, et le renvoi
%  reste du texte ordinaire. C'est un des trois ecarts declares dans
%  outils/renvoji.py.
% ===================================================================

%%K t06-apar-1 apar
\VUsaut{6.0mm}
\VUtitre{15.0pt}{60}{TABELO N\textsuperscript{o} 6}
\VUsaut{1.4mm}
\VUtitre{11.4pt}{0}{La interna domo, la meblaro, la nutrado,}
\VUsaut{0.4mm}
\VUtitre{11.4pt}{0}{la lumigado.}
\VUsaut{2.2mm}

%%K t06-apar-2 apar
La domo estas finkonstruita. La onklo de Johano estas instalita en sia nova
loĝejo, kies aranĝon ni konas. Nun ni vidu kiel li meblis sian domon. Ni
esploros unue:

%%K t06-tit-1 sub
\VUpk{10.2pt}{40}{La dormoĉambro.}
\VUsaut{3.0mm}

%%K t06-c1-01-1 p
1. --- Ni alvenas maloportune, ĉar la ĉambristino estas okupata purigi la
\VUgras{ĉambron}.

%%K t06-c1-02-1 p
2. --- Sur la \VUgras{tualettablo} \textsuperscript{(1)} estas tie kaj tie
diversaj objektoj per kiuj oni lavas sin, kombas sin, mallonge oni tualetas.
Ili estas la \VUgras{pelveto} \textsuperscript{(2)} kaj la \VUgras{akvokruĉo}
\textsuperscript{(3)}, la \VUgras{kapbroso} \textsuperscript{(4)}, la
\VUgras{kombilo} \textsuperscript{(5)}, la \VUgras{sapo} \textsuperscript{(6)}
kaj la \VUgras{spongo} \textsuperscript{(7)}, fine \VUgras{flakoneto}
\textsuperscript{(8)} por la dentifrico likva.

%%K t06-c1-03-1 p
3. --- Sur la planko, antaŭ la tualettablo, jen la \VUgras{sitelo}
\textsuperscript{(9)} por kolekti la malpurajn akvojn.

%%K t06-c1-04-1 p
4. --- Ni vidas en la \VUgras{antaŭĉambro} \textsuperscript{(10)} kelkajn
\VUgras{vesthokojn} \textsuperscript{(13)} kaj du \VUgras{ombrelojn}
\textsuperscript{(11)} en la \VUgras{ombrelujo} \textsuperscript{(12)}.

%%K t06-c1-05-1 p
5. --- Aliflanke de la pordo estas la \VUgras{spegulŝranko}
\textsuperscript{(14)} kiu enhavas la \VUgras{tolaĵaron}
\textsuperscript{(15)}.

%%K t06-c1-06-1 p
6. --- Ni profitu ke la \VUgras{ĉambristino} \textsuperscript{(16)} aranĝas la
\VUgras{liton} \textsuperscript{(17)}, por esplori ĝiajn diversajn partojn. La
\VUgras{litframo} \textsuperscript{(20)} enhavas bonan \VUgras{risortmatracon}
\textsuperscript{(21)} sur kiun Justina tuj metos lanan \VUgras{matracon}
\textsuperscript{(22)}. Poste ŝi prenos de la seĝoj la du \VUgras{littukojn}
\textsuperscript{(23)} kaj la \VUgras{kovrilon} \textsuperscript{(24)}. Tiam ŝi
metos sur la litkapon, la \VUgras{transversan kusenon} \textsuperscript{(25)},
la \VUgras{kapkusenon} \textsuperscript{(26)} kiuj estas kun la
\VUgras{lanugkovrilo} \textsuperscript{(27)} sur la \VUgras{tapiŝo lita}
\textsuperscript{(32)}. Ŝi uzas plumfaskon por senpolvigi la \VUgras{kurtenojn}
\textsuperscript{(19)} kaj la \VUgras{litbaldakenon} \textsuperscript{(18)} kaj
jen parto de la laboro finita.

%%K t06-c1-07-1 p
7. --- Tiam ŝi devos ordigi iom la \VUgras{littableton}
\textsuperscript{(28)}, restreĉi la \VUgras{vekhorloĝon} \textsuperscript{(29)},
prepari la \VUgras{noktolampeton} \textsuperscript{(30)}, depreni la
\VUgras{kandelon} \textsuperscript{(31)} por purigi la kandelingon.

%%K t06-tit-2 sub
\VUsaut{5.0mm}
\VUpk{10.2pt}{40}{La laborkorbo kaj la kamenaj akcesoraĵoj.}
\VUsaut{3.0mm}

%%K t06-c1-08-1 p
8. --- La \VUgras{laborkorbo} \textsuperscript{(34)} apud la \VUgras{tabureto}
\textsuperscript{(33)} ankaŭ tre bezonas esti ordigita. Ŝi devos faldi la
\VUgras{brodaĵojn} \textsuperscript{(35)}, enfermi en la \VUgras{labortablo}
\textsuperscript{(38)} la \VUgras{fingringon}, la \VUgras{fadenon}
\textsuperscript{(36)}, la \VUgras{kudrilojn} \textsuperscript{(37)} kaj la
\VUgras{tondilon} \textsuperscript{(39)} kaj fermi per la \VUgras{ŝlosilo}
\textsuperscript{(40)} la kovrilon de la tablo.

%%K t06-c1-09-1 p
9. --- Sur la \VUgras{unupieda tablo} \textsuperscript{(41)} de la mezo,
Justina renovigos la akvon de la \VUgras{karafo} \textsuperscript{(42)}. Ŝi
metos \VUgras{sukeron} en la \VUgras{sukerujon} \textsuperscript{(43)},
purigos la \VUgras{sukerprenilon} \textsuperscript{(44)} kaj deprenos la
\VUgras{vestbroson} \textsuperscript{(46)}.

%%K t06-c1-10-1 p
10. --- La dekstra flanko de la ĉambro postulos de ŝi malpli multe da laboro,
ĉar la \VUgras{longa brakseĝo} \textsuperscript{(47)} kaj ĝia \VUgras{kuseno}
\textsuperscript{(48)} estas en sia ĝusta loko, kaj, ĉar ne estas malvarma
vetero, nenio estas dislokita inter la akcesoraĵoj (ilaroj) de la
\VUgras{kameno} \textsuperscript{(49)}. \VUgras{Ŝovelilo}
\textsuperscript{(50)}, \VUgras{prenegiloj} \textsuperscript{(51)},
\VUgras{ŝtiptenilo} \textsuperscript{(55)}, \VUgras{cindrobarilo}
\textsuperscript{(56)} estas puraj; la \VUgras{fajroŝirmilo}
\textsuperscript{(54)} ne estis dislokita kaj la \VUgras{lignokesto}
\textsuperscript{(52)} estas bone provizita per \VUgras{brulŝtipoj}
\textsuperscript{(53)}.

%%K t06-noto-1 noto
\VUnotes{143.2mm}{%
(*) Babo = infaneto, la A. baby, F. bébé, I. bambino.%
}

%%K t06-noto-2 noto
\VUnotes{143.2mm}{%
(*) Lambrequino = F. lambrequin, H. lambrequines, Port. lambrequins, eĉ G. A.
lambrequins, do G. A. F. P. H.%
}

%%K t06-c1-11-1 p
11. --- Tamen ŝi devos senpolvigi la \VUgras{pendolhorloĝon}
\textsuperscript{(57)}, la \VUgras{kandelabrojn} \textsuperscript{(58)} kaj la
\VUgras{poŝaĵujojn} \textsuperscript{(59)}, fine viŝi la \VUgras{spegulon}
\textsuperscript{(60)}.

%%K t06-c1-12-1 p
12. --- Pri la \VUgras{lulilo} \textsuperscript{(61)} de la infaneto (de la
babo (*)), ĝi jam estas preta.

%%K t06-tit-3 sub
\VUsaut{5.0mm}
\VUpk{10.2pt}{40}{La salono.}
\VUsaut{3.0mm}

%%K t06-c2-01-1 p
1. --- Nun jen la salono. Ĝi estas tre bela \VUgras{ĉambro}. La kameno, kiu
situas ĉe la dekstra angulo estas ornamita per delikata \VUgras{statueto}
\textsuperscript{(1)}, per du belaj \VUgras{vazoj} \textsuperscript{(3)} kaj
drapirita per velura \VUgras{lambrekeno} \textsuperscript{(2)} (*).

%%K t06-c2-02-1 p
2. --- Por malhelpi ke la fajreroj de la fajrujo ekbruligu la tapiŝon, oni
metis antaŭ la fajrujon kupran \VUgras{fajreroŝirmilon} \textsuperscript{(4)}.

%%K t06-c2-03-1 p
3. --- Proksime, eleganta \VUgras{skribotablo} \textsuperscript{(5)} sinjorina
estas malfermita. Sur ĝi la \VUgras{domomastrino} \textsuperscript{(23)}
skribas siajn leterojn.

%%K t06-c2-04-1 p
4. --- La fenestro estas garnita per \VUgras{vitraj kurtenoj}
\textsuperscript{(6)} kun \VUgras{ligiloj} \textsuperscript{(8)} alkroĉitaj al
\VUgras{hokoj} \textsuperscript{(7)} kaj per ŝtofaj kurtenegoj \VUgras{supre
drapiritaj} \textsuperscript{(9)}.

%%K t06-c2-05-1 p
5. --- En la embrazuro de la fenestro, \VUgras{potujo} \textsuperscript{(11)}
enhavas verdan \VUgras{planton} \textsuperscript{(10)}, belegan palmarbon kies
folioj sin etendas preskaŭ ĝis la \VUgras{petrollampo} \textsuperscript{(12)}.

%%K t06-c2-06-1 p
6. --- La \VUgras{lumŝirmilo} \textsuperscript{(13)} de la lampo estas aranĝita
de Maria, la ĉarma \VUgras{junulino} \textsuperscript{(14)} kiu estas ĉe la
piano \textsuperscript{(15)}. Tre rekte sidanta sur la \VUgras{tabureto}
\textsuperscript{(16)}, ŝi studas muzikaĵon. Ŝi estas tre lerta kaj zorgema,
kiel pruvas ŝia \VUgras{muzikŝranko} \textsuperscript{(17)} kiu estas perfekte
ordigita.

%%K t06-tit-4 sub
\VUsaut{5.0mm}
\VUpk{10.2pt}{40}{La bubaĉo.}
\VUsaut{3.0mm}

%%K t06-c2-07-1 p
7. --- Ŝia frato, Paŭlo, serĉas \VUgras{libron} \textsuperscript{(19)} en la
\VUgras{biblioteko} \textsuperscript{(18)}. Li estas \VUgras{junulo}
\textsuperscript{(20)}, movema kaj netolerebla. Alkroĉante sin al la
biblioteko por atingi la libron kiu estas tro supre, li riskas faligi la
\VUgras{buston} \textsuperscript{(21)} kiu estas supre.

%%K t06-c2-08-1 p
8. --- Li profitas, por fari tiun stultaĵon, ke lia patrino estas okupata
konversacii kun amikino, \VUgras{sinjorino} \textsuperscript{(22)} kiu venis
pasigi kelkajn tagojn en ŝia hejmo. La patrino sidas sur \VUgras{kanapo}
\textsuperscript{(24)} apud la \VUgras{brakseĝo} \textsuperscript{(25)}, kaj
ŝia amikino estas sur la \VUgras{seĝo} \textsuperscript{(27)} de kiu oni vidas
nur la \VUgras{dorsapogilon} \textsuperscript{(26)}.

%%K t06-c2-09-1 p
9. --- La granda \VUgras{kadro} \textsuperscript{(30)} pendanta ĉe la muro
enhavas la \VUgras{portreton} \textsuperscript{(29)} de parencino. En la
\VUgras{albumo} \textsuperscript{(32)} metita sur la tablo estas kunigitaj la
\VUgras{fotografaĵoj} \textsuperscript{(33)} de la ceteraj familianoj. La
infanoj ĵus foliumis ĝin, ĉar la \VUgras{seĝoj} \textsuperscript{(31)} estas
ankoraŭ grupigitaj ĉirkaŭ la tablo.

%%K t06-c2-10-1 p
10. --- La \VUgras{ĉina vazo} \textsuperscript{(34)} porcelana, kiu estas apud
la \VUgras{papertranĉilo} \textsuperscript{(35)} estis donacita al Maria por
ŝia festo, kaj la \VUgras{paravento} \textsuperscript{(36)} kiu garnas la
angulon de la ĉambro estas alportita el Ĉinio de ŝia onklo.

%%K t06-c2-11-1 p
11. --- Gajigata de la belegaj \VUgras{pentraĵoj} \textsuperscript{(37)}
alkroĉitaj ĉe la \VUgras{vando} \textsuperscript{(42)}, lumigata de la
\VUgras{plafonlustro} \textsuperscript{(38)}, kies \VUgras{elektraj lampoj}
\textsuperscript{(39)} vespere brilas ĝoje, tiu salono aspektas tre agrabla. La
mola \VUgras{tapiŝo} \textsuperscript{(40)} etendita sur la parketo, kaj la tiel
facile uzebla \VUgras{elektra sonorilo} \textsuperscript{(41)}, finfaras ĝian
komforton.

%%K t06-tit-5 sub
\VUsaut{3.6mm}
\VUpk{10.2pt}{40}{La manĝoĉambro.}
\VUsaut{3.0mm}

%%K t06-c3-01-1 p
1. --- Sonis la vespermanĝa horo. La tuta familio, escepte Maria, estas
kunigita ĉirkaŭ la tablo. La manĝoĉambro estas agrable varmigita de la bonega
\VUgras{forno} \textsuperscript{(1)} fajenca, kun maldika \VUgras{tubo}
\textsuperscript{(2)}.

%%K t06-c3-02-1 p
2. --- Tiu ĉambro enhavas ne multajn meblojn. Laŭlonge de la muro pendas,
super la \VUgras{benketo} \textsuperscript{(4)}, \VUgras{fonteto}
\textsuperscript{(3)} ornama. Tre proksime, estas la \VUgras{kredenco}
\textsuperscript{(5)} sur kiun oni metas la malvarmajn manĝaĵojn, la
desertpladojn, kaj la diversajn tablilarojn kiujn oni ne tuj bezonas.

%%K t06-c3-03-1 p
3. --- Tiel ni vidas sur la supra breto, \VUgras{kokidon rostitan}
\textsuperscript{(7)}, \VUgras{fiŝon} \textsuperscript{(8)} kaj \VUgras{pelvon}
\textsuperscript{(10)} kun \VUgras{fruktoj} \textsuperscript{(9)}. Sube jen la
\VUgras{fromaĝo} \textsuperscript{(11)}, la \VUgras{pano}
\textsuperscript{(12)}, la \VUgras{flakonujo} \textsuperscript{(13)} kiu
enhavas ĉion bezonatan por spici la salaton: la oleo, la vinagro, la
\VUgras{salo} \textsuperscript{(14)} kaj la \VUgras{pipro}
\textsuperscript{(15)}. Fine, sur la malsupra breto estas \VUgras{kuko}
\textsuperscript{(17)} apud la \VUgras{mustardujo} \textsuperscript{(16)}.

%%K t06-c3-04-1 p
4. --- La manĝoĉambra horloĝo, kiun oni nomas \VUgras{murhorloĝo}
\textsuperscript{(20)}, havas tre apartan formon. Ĝi estas enkadrigita en ligna
kadro kiu krome enhavas \VUgras{barometron} \textsuperscript{(19)} kaj
\VUgras{termometron} \textsuperscript{(18)}.

%%K t06-tit-6 sub
\VUsaut{4.6mm}
\VUpk{10.2pt}{40}{La vespermanĝa servo.}
\VUsaut{3.0mm}

%%K t06-c3-05-1 p
5. --- Ĉe ĉiuj muroj de la ĉambro estas aplikitaj \VUgras{lignaj paneloj}
\textsuperscript{(21)} el vaksita kverko. Ŝtofa \VUgras{pordotapeto}
\textsuperscript{(22)} sufiĉe fermas la pordon kiu donas eniron al la verando.
Tie la familio iros trinki la kafon post la vespermanĝo. Jam la \VUgras{tasoj}
\textsuperscript{(24)} kaj la \VUgras{subtasoj} \textsuperscript{(25)} estas
pretaj sur la \VUgras{pladmeblo} \textsuperscript{(23)}.

%%K t06-c3-06-1 p
6. --- Super la tablo estas fiksita al la plafono \VUgras{pendolumigilo}
\textsuperscript{(26)} kies tre brilanta lampo vespere lumigas la tutan
manĝoĉambron.

%%K t06-c3-07-1 p
7. --- Nun ni esploru la \VUgras{manĝotablon} \textsuperscript{(27)} mem. Kiam
eniris la familio, la manĝilaro estis preta. Sur la \VUgras{tablotuko}
\textsuperscript{(28)} estis metita, ĉe la loko de ĉiu kunmanĝonto,
\VUgras{telero} \textsuperscript{(33)}, \VUgras{buŝtuko}
\textsuperscript{(29)}, \VUgras{kulero} \textsuperscript{(34)},
\VUgras{forketo} \textsuperscript{(35)}, \VUgras{tranĉilo}
\textsuperscript{(36)} kaj \VUgras{glaso} \textsuperscript{(32)}. Estis ankaŭ
\VUgras{boteloj} \textsuperscript{(30)} por vino kaj \VUgras{karafo}
\textsuperscript{(31)} por akvo.

%%K t06-c3-08-1 p
8. --- La \VUgras{servisto} \textsuperscript{(39)} unue alportis la
\VUgras{supujon} \textsuperscript{(37)}, en kiu ankoraŭ fumas iom da supo. Nun
li servas la unuan \VUgras{pladon} \textsuperscript{(38)}, viandopladon. Poste
li prezentos tiujn kiujn ni jam nomis, fine, post la \VUgras{desertoplado}, li
profitos ke liaj mastroj estos enirintaj la verandon, por depreni la
tablilarojn kaj reordigi la manĝoĉambron.

%%K t06-c3-08-2 p
\textit{Noto.} --- \textit{La komunuzaj vortoj rilatantaj la nutraĵaron, ĉi
tie estas indikitaj tre resume. La listo estos kompletigata sur la du tabeloj
« La Hotelo » n\textsuperscript{o} 13) kaj « La Bazaro » (n\textsuperscript{o}
15).}

%%K t06-tit-7 sub
\VUsaut{4.6mm}
\VUpk{10.2pt}{40}{La kuirejo.}
\VUsaut{3.0mm}

%%K t06-c4-01-1 p
1. --- En la kuirejo, kiun ni ekrigardas tra la pordo duonmalfermita, ni vidas
pentrindan malordon. Antaŭ la \VUgras{fajrujo} \textsuperscript{(1)} en kiu
\VUgras{fajro} \textsuperscript{(3)} krakas, estas instalita
\VUgras{rostoturnilo} \textsuperscript{(5)}. Bela \VUgras{rostaĵo}
\textsuperscript{(7)} estas \VUgras{surrostostangigita} \textsuperscript{(6)}
kaj finkuiriĝas.

%%K t06-c4-02-1 p
2. --- La \VUgras{blovilo} \textsuperscript{(8)} kaj la \VUgras{karboŝovelilo}
\textsuperscript{(9)}, kiujn oni ĵus uzis, restis sur la planko same kiel la
\VUgras{brulŝtipoj} \textsuperscript{(10)}.

%%K t06-c4-03-1 p
3. --- La supro de la kameno estas ordigita; la \VUgras{lanterno}
\textsuperscript{(11)}, la \VUgras{alumetujo} \textsuperscript{(13)}, en kiu
estas la \VUgras{alumetoj} \textsuperscript{(12)}, la \VUgras{kandelingo}
\textsuperscript{(14)} kaj la \VUgras{mankandelingo} \textsuperscript{(15)} kun
\VUgras{kandelo} \textsuperscript{(16)} estas en sia loko. Sed la granda
\VUgras{karbositelo} \textsuperscript{(18)}, plena de \VUgras{terkarbo}
\textsuperscript{(17)} kiun la \VUgras{kuiristino} \textsuperscript{(23)} ĵus
plenigis en la kelo, restis meze de la kuirejo.

%%K t06-c4-04-1 p
4. --- Vere, la kompatinda virino devas rapidi por ke la tagmanĝo estu preta
ĝustatempe. Nun ŝi estas apud la \VUgras{kuirforno} \textsuperscript{(19)} kaj
kirlas \VUgras{saŭcon} \textsuperscript{(24)} kiun konvenas ne bruligi tro.

%%K t06-c4-05-1 p
5. --- En la \VUgras{bakforno} \textsuperscript{(20)} bakiĝas kuko, kiun ŝi
preparis por la manĝeto de la infanoj. Super la kuirforno pendas la
\VUgras{potkulero} \textsuperscript{(21)} kaj la \VUgras{ŝaŭmkulero}
\textsuperscript{(22)}.

%%K t06-tit-8 sub
\VUsaut{4.0mm}
\VUpk{10.2pt}{40}{La vazaro.}
\VUsaut{3.0mm}

%%K t06-c4-06-1 p
6. --- La \VUgras{kuirilaro} \textsuperscript{(25)}, alkroĉita en la fundo de
la ĉambro, estas tre pura. La belaj kupraj \VUgras{kaseroloj}
\textsuperscript{(26)}, la \VUgras{laktopoto} \textsuperscript{(27)}, la
\VUgras{kribrileto} \textsuperscript{(28)} kaj la \VUgras{raspilo}
\textsuperscript{(29)} estis zorge purigitaj.

%%K t06-c4-07-1 p
7. --- La \VUgras{salujo} \textsuperscript{(30)} estas metita sur la
pladŝrankon apud la \VUgras{tekruĉo} \textsuperscript{(31)}, kaj la
\VUgras{oleobotelego} \textsuperscript{(32)}. La \VUgras{tirkesto}
\textsuperscript{(33)} kredeble enhavas la manĝilarojn kaj la tranĉilojn de
kuirejo.

%%K t06-c4-08-1 p
8. --- La \VUgras{balailo} \textsuperscript{(34)} kaj la \VUgras{plumfasko}
\textsuperscript{(35)} estas en angulo. La kuiristino ĵus uzis la
\VUgras{lignoblokon} \textsuperscript{(36)}, ŝi lasis ĝin apud la fenestro.

%%K t06-c4-09-1 p
9. --- \VUgras{Manviŝilo} \textsuperscript{(37)} pendas ĉe la muro, apud la
\VUgras{funelo} \textsuperscript{(38)}. En tiu parto de la kuirejo estas
formetitaj la \VUgras{kradrostilo} \textsuperscript{(39)}, la
\VUgras{haktranĉilo} \textsuperscript{(40)} kaj la \VUgras{hakplanko}
\textsuperscript{(41)}. Poste, super la haktranĉilo, sur \VUgras{breto}
\textsuperscript{(42)}, estas \VUgras{potoj} \textsuperscript{(43)}, la
\VUgras{bolpoto} \textsuperscript{(44)} kun sia \VUgras{kovrilo}
\textsuperscript{(45)} kaj la \VUgras{marmito} \textsuperscript{(46)}. La
\VUgras{pato} \textsuperscript{(47)} pendas proksime.

%%K t06-c4-10-1 p
10. --- Nur la kupra \VUgras{krano} \textsuperscript{(48)} ĵetas brilon en tiu
ĉi angulo. La \VUgras{lavujo} \textsuperscript{(49)} sur kiu estas metita la
dika \VUgras{kruĉego} \textsuperscript{(50)}, estas kredeble tre komforta kun
sia malgranda ŝranko, en kiu oni konservas la \VUgras{vinagrobarelon}
\textsuperscript{(51)} provizitan per sia \VUgras{krano}
\textsuperscript{(52)}, la \VUgras{restaĵujon} \textsuperscript{(53)} kaj la
\VUgras{malpurajn teleroj} \textsuperscript{(54)}.

%%K t06-tit-9 sub
\VUsaut{4.0mm}
\VUpk{10.2pt}{40}{Aliaj objektoj metitaj en la kuirejo.}
\VUsaut{3.0mm}

%%K t06-c4-11-1 p
11. --- La \VUgras{kuvego} \textsuperscript{(55)} en kiu oni lavas la
\VUgras{legomojn} \textsuperscript{(58)} kaj la gisfera \VUgras{kaldrono}
\textsuperscript{(56)} metita sur la \VUgras{kahelplankon}
\textsuperscript{(57)} kredeble havas ankaŭ sian lokon en tiu ŝranko.

%%K t06-c4-12-1 p
12. --- La kuiristino certe ne malhavas fornojn, ĉar jen krome, sur la angulo
de la tablo, \VUgras{gasforno} \textsuperscript{(59)} sur kiu varmiĝas akvo en
malgranda kasrolo. La \VUgras{vaporo} \textsuperscript{(60)} kiu eskapas el ĝi
montras al ni, ke kredeble ĝi bolas. Verŝajne ĉi tiu akvo estos uzata por la
kafo, ĉar jen sur la alia ekstremo de la tablo, la \VUgras{kafkruĉo}
\textsuperscript{(64)} kaj la \VUgras{kafmuelilo} \textsuperscript{(63)}.

%%K t06-c4-13-1 p
13. --- La forno estas konektita kun la \VUgras{gasbrulilo}
\textsuperscript{(62)} per longa \VUgras{kaŭĉuka tubeto}
\textsuperscript{(61)}.

%%K t06-c4-14-1 p
14. --- La \VUgras{taga servistino} \textsuperscript{(66)} kiu venas helpi la
kuiristinon, preparas la legomojn por la vespermanĝo. Kiam ili estos
senŝeligitaj kaj lavitaj, ŝi metos ilin en la \VUgras{basenon}
\textsuperscript{(65)} atendante la horon por kuiri ilin.

%%K t06-tit-10 sub
\VUsaut{4.0mm}
{\centering\fontsize{10.2pt}{10.2pt}\selectfont\textls[40]{\scshape La banejo.} \textit{(Banĉambro.)}\par}
\VUsaut{3.0mm}

%%K t06-c5-01-1 p
1. --- Al ni restas farenda nur unu maldiskretaĵo por ekkoni la ĉefajn
ĉambrojn de la apartamento: vidi la instalaĵon de la banejo. Tio ne postulos
longan tempon, ĉar ĝi ne estas granda, kaj ni rapide scios ke la
\VUgras{bankuvo} \textsuperscript{(1)} havas bonan \VUgras{banvarmigilon}
\textsuperscript{(2)}, ke la \VUgras{sapkonko} \textsuperscript{(3)} estas
atingebla de la mano de banulo kaj ke la \VUgras{tukportilo}
\textsuperscript{(5)} certe igos facile sekigi la \VUgras{tualettukojn}
\textsuperscript{(4)}.

%%K t06-c5-02-1 p
2. --- Tiu ĉambro havas siajn murojn kovritajn per \VUgras{fajencaĵoj}
\textsuperscript{(6)}, kiuj ebligis instali \VUgras{duŝaparaton}
\textsuperscript{(7)} sen iel timi ke la akvo, kiu respruĉas dum ĝia uzado,
difektos la murojn. Zinka \VUgras{baseneto} \textsuperscript{(8)} kiu servas
por la piedbanoj, kompletigas la meblaron.
\VUsaut{6.0mm}
\VUfilet{22mm}
